% Options for packages loaded elsewhere
\PassOptionsToPackage{unicode}{hyperref}
\PassOptionsToPackage{hyphens}{url}
%
\documentclass[
  11pt,
  a5paper,
]{book}

\usepackage{amsmath,amssymb}
\usepackage{setspace}
\usepackage{iftex}
\ifPDFTeX
  \usepackage[T1]{fontenc}
  \usepackage[utf8]{inputenc}
  \usepackage{textcomp} % provide euro and other symbols
\else % if luatex or xetex
  \usepackage{unicode-math}
  \defaultfontfeatures{Scale=MatchLowercase}
  \defaultfontfeatures[\rmfamily]{Ligatures=TeX,Scale=1}
\fi
\usepackage{lmodern}
\ifPDFTeX\else  
    % xetex/luatex font selection
    \setmainfont[]{EB Garamond}
\fi
% Use upquote if available, for straight quotes in verbatim environments
\IfFileExists{upquote.sty}{\usepackage{upquote}}{}
\IfFileExists{microtype.sty}{% use microtype if available
  \usepackage[]{microtype}
  \UseMicrotypeSet[protrusion]{basicmath} % disable protrusion for tt fonts
}{}
\makeatletter
\@ifundefined{KOMAClassName}{% if non-KOMA class
  \IfFileExists{parskip.sty}{%
    \usepackage{parskip}
  }{% else
    \setlength{\parindent}{0pt}
    \setlength{\parskip}{6pt plus 2pt minus 1pt}}
}{% if KOMA class
  \KOMAoptions{parskip=half}}
\makeatother
\usepackage{xcolor}
\usepackage[lmargin=2.5cm,rmargin=2.5cm,tmargin=3cm,bmargin=3cm]{geometry}
\setlength{\emergencystretch}{3em} % prevent overfull lines
\setcounter{secnumdepth}{-\maxdimen} % remove section numbering
% Make \paragraph and \subparagraph free-standing
\makeatletter
\ifx\paragraph\undefined\else
  \let\oldparagraph\paragraph
  \renewcommand{\paragraph}{
    \@ifstar
      \xxxParagraphStar
      \xxxParagraphNoStar
  }
  \newcommand{\xxxParagraphStar}[1]{\oldparagraph*{#1}\mbox{}}
  \newcommand{\xxxParagraphNoStar}[1]{\oldparagraph{#1}\mbox{}}
\fi
\ifx\subparagraph\undefined\else
  \let\oldsubparagraph\subparagraph
  \renewcommand{\subparagraph}{
    \@ifstar
      \xxxSubParagraphStar
      \xxxSubParagraphNoStar
  }
  \newcommand{\xxxSubParagraphStar}[1]{\oldsubparagraph*{#1}\mbox{}}
  \newcommand{\xxxSubParagraphNoStar}[1]{\oldsubparagraph{#1}\mbox{}}
\fi
\makeatother


\providecommand{\tightlist}{%
  \setlength{\itemsep}{0pt}\setlength{\parskip}{0pt}}\usepackage{longtable,booktabs,array}
\usepackage{calc} % for calculating minipage widths
% Correct order of tables after \paragraph or \subparagraph
\usepackage{etoolbox}
\makeatletter
\patchcmd\longtable{\par}{\if@noskipsec\mbox{}\fi\par}{}{}
\makeatother
% Allow footnotes in longtable head/foot
\IfFileExists{footnotehyper.sty}{\usepackage{footnotehyper}}{\usepackage{footnote}}
\makesavenoteenv{longtable}
\usepackage{graphicx}
\makeatletter
\newsavebox\pandoc@box
\newcommand*\pandocbounded[1]{% scales image to fit in text height/width
  \sbox\pandoc@box{#1}%
  \Gscale@div\@tempa{\textheight}{\dimexpr\ht\pandoc@box+\dp\pandoc@box\relax}%
  \Gscale@div\@tempb{\linewidth}{\wd\pandoc@box}%
  \ifdim\@tempb\p@<\@tempa\p@\let\@tempa\@tempb\fi% select the smaller of both
  \ifdim\@tempa\p@<\p@\scalebox{\@tempa}{\usebox\pandoc@box}%
  \else\usebox{\pandoc@box}%
  \fi%
}
% Set default figure placement to htbp
\def\fps@figure{htbp}
\makeatother

\makeatletter
\@ifpackageloaded{bookmark}{}{\usepackage{bookmark}}
\makeatother
\makeatletter
\@ifpackageloaded{caption}{}{\usepackage{caption}}
\AtBeginDocument{%
\ifdefined\contentsname
  \renewcommand*\contentsname{Índice}
\else
  \newcommand\contentsname{Índice}
\fi
\ifdefined\listfigurename
  \renewcommand*\listfigurename{Lista de Figuras}
\else
  \newcommand\listfigurename{Lista de Figuras}
\fi
\ifdefined\listtablename
  \renewcommand*\listtablename{Lista de Tabelas}
\else
  \newcommand\listtablename{Lista de Tabelas}
\fi
\ifdefined\figurename
  \renewcommand*\figurename{Figura}
\else
  \newcommand\figurename{Figura}
\fi
\ifdefined\tablename
  \renewcommand*\tablename{Tabela}
\else
  \newcommand\tablename{Tabela}
\fi
}
\@ifpackageloaded{float}{}{\usepackage{float}}
\floatstyle{ruled}
\@ifundefined{c@chapter}{\newfloat{codelisting}{h}{lop}}{\newfloat{codelisting}{h}{lop}[chapter]}
\floatname{codelisting}{Listagem}
\newcommand*\listoflistings{\listof{codelisting}{Lista de Listagens}}
\makeatother
\makeatletter
\makeatother
\makeatletter
\@ifpackageloaded{caption}{}{\usepackage{caption}}
\@ifpackageloaded{subcaption}{}{\usepackage{subcaption}}
\makeatother

\ifLuaTeX
\usepackage[bidi=basic]{babel}
\else
\usepackage[bidi=default]{babel}
\fi
\babelprovide[main,import]{portuguese}
\ifPDFTeX
\else
\babelfont{rm}[]{EB Garamond}
\fi
% get rid of language-specific shorthands (see #6817):
\let\LanguageShortHands\languageshorthands
\def\languageshorthands#1{}
\usepackage{bookmark}

\IfFileExists{xurl.sty}{\usepackage{xurl}}{} % add URL line breaks if available
\urlstyle{same} % disable monospaced font for URLs
\hypersetup{
  pdftitle={Brain Drafts},
  pdfauthor={bmhash},
  pdflang={pt},
  hidelinks,
  pdfcreator={LaTeX via pandoc}}


\title{Brain Drafts}
\usepackage{etoolbox}
\makeatletter
\providecommand{\subtitle}[1]{% add subtitle to \maketitle
  \apptocmd{\@title}{\par {\large #1 \par}}{}{}
}
\makeatother
\subtitle{Notas para a minha filha}
\author{bmhash}
\date{16 maio 2026}

\begin{document}
\frontmatter
\maketitle

\renewcommand*\contentsname{Índice}
{
\setcounter{tocdepth}{2}
\tableofcontents
}

\setstretch{1.5}
\mainmatter
\bookmarksetup{startatroot}

\chapter{Para Ti}\label{para-ti}

\emph{«O processo natural de transmissão de geração em geração assegura
a perpetuação das características físicas\ldots{} mas a educação é
elevada a uma potência muito superior pelo esforço deliberado do
conhecimento e da vontade humana para atingir um fim conhecido.»}

--- Werner Jaeger, \emph{Paideia}

Este livro não é um manual. Não é um conjunto de regras. Não é um
testamento.

É uma tentativa --- imperfeita, inacabada, honesta --- de pôr em
palavras aquilo que sei ser verdade. Não porque o mundo mo disse, mas
porque o vivi, o pensei, e decidi que merecia ser transmitido.

Marco Aurélio escreveu as suas \emph{Meditações} para si mesmo --- notas
privadas sobre como viver, nunca destinadas ao mundo. Eu escrevo para
ti. Com a mesma qualidade de honestidade, de directeza, e de ausência de
performance.

Não espero que concordes com tudo. Espero que, num momento de confusão
ou de dor, alguma destas linhas te lembre de quem és quando ninguém está
a ver.

O livro está organizado em três círculos --- do eu interior para os
outros, dos outros para o mundo:

\begin{itemize}
\tightlist
\item
  \textbf{Logos} --- quem és quando estás sozinha com os teus
  pensamentos
\item
  \textbf{Ethos} --- como te encontras com os outros e com o mundo
\item
  \textbf{Cosmos} --- como te situas num universo maior do que tu
\end{itemize}

Não há pressa. Lê quando quiseres. Relê quando precisares. E se um dia
achares que isto já não te serve --- isso também está certo. O melhor
que este livro pode fazer é, eventualmente, tornar-se desnecessário.

\part{Círculo I --- O Eu Interior (Logos)}

\chapter{Não és o que te acontece}\label{nuxe3o-uxe9s-o-que-te-acontece}

\emph{«Não são os acontecimentos que perturbam as pessoas; são os
julgamentos que fazem sobre eles.»}

--- Marco Aurélio

\textbf{Âncora Estóica:} Dicotomia do Controlo\\
\textbf{Conceito Grego:} \emph{Prohairesis} (escolha moral, vontade)

\chapter{A opinião que formas é a
ferida}\label{a-opiniuxe3o-que-formas-uxe9-a-ferida}

\emph{«Retira a tua opinião, e retiras a queixa `fui ferido'. Retira a
queixa `fui ferido', e a ferida desaparece.»}

--- Marco Aurélio

\textbf{Âncora Estóica:} Disciplina do Assentimento\\
\textbf{Conceito Grego:} \emph{Logos} (razão, palavra interior)

\chapter{Sobre o que é verdadeiramente
bom}\label{sobre-o-que-uxe9-verdadeiramente-bom}

\emph{«A virtude é o único bem verdadeiro.»}

--- Zenão de Cítio

\textbf{Âncora Estóica:} A virtude como único bem\\
\textbf{Conceito Grego:} \emph{Eudaimonia} (florescimento humano)

\chapter{A admiração é uma forma de
inteligência}\label{a-admirauxe7uxe3o-uxe9-uma-forma-de-inteliguxeancia}

\emph{«A filosofia começa na admiração.»}

--- Platão

\textbf{Âncora Estóica:} A vida examinada\\
\textbf{Conceito Grego:} \emph{Thaumazein} (admiração, espanto)

\part{Círculo II --- Com os Outros (Ethos)}

\chapter{A coragem não é a ausência do
medo}\label{a-coragem-nuxe3o-uxe9-a-ausuxeancia-do-medo}

\emph{«De que vale a coragem se não está ao serviço da justiça?»}

--- Séneca

\textbf{Âncora Estóica:} A virtude da \emph{Andreia} (coragem)\\
\textbf{Conceito Grego:} \emph{Areté} (excelência, virtude)

\chapter{Como tratar um estranho}\label{como-tratar-um-estranho}

\emph{«O que não é bom para a colmeia não pode ser bom para a abelha.»}

--- Marco Aurélio

\textbf{Âncora Estóica:} Justiça e o bem comum\\
\textbf{Conceito Grego:} \emph{Dikaiosynê} (justiça)

\chapter{Suficiente}\label{suficiente}

\emph{«A riqueza consiste não em ter grandes posses, mas em ter poucas
necessidades.»}

--- Epicteto

\textbf{Âncora Estóica:} Temperança e moderação\\
\textbf{Conceito Grego:} \emph{Sôphrosynê} (autodomínio, temperança)

\chapter{Sobre o amor e as suas
exigências}\label{sobre-o-amor-e-as-suas-exiguxeancias}

\emph{«O amor não é apenas um sentimento. É um acto de vontade.»}

\textbf{Âncora Estóica:} O amor como sacrifício\\
\textbf{Conceito Grego:} \emph{Philia} / \emph{Eros} (amizade profunda /
amor)

\chapter{Sobre o corpo}\label{sobre-o-corpo}

\emph{«Nenhuma alma pode alcançar a excelência sem um corpo saudável.»}

--- Platão

\textbf{Âncora Estóica:} \emph{Areté} física, a saúde como disciplina\\
\textbf{Conceito Grego:} \emph{Sophrosyne} (moderação, equilíbrio)

\chapter{O que deves à beleza}\label{o-que-deves-uxe0-beleza}

\emph{«A beleza é o esplendor da verdade.»}

--- Platão

\textbf{Âncora Estóica:} Arte, música, admiração no mundo\\
\textbf{Conceito Grego:} \emph{Kallos} (beleza)

\chapter{A comunidade a que
pertences}\label{a-comunidade-a-que-pertences}

\emph{«O homem é por natureza um animal político.»}

--- Aristóteles

\textbf{Âncora Estóica:} \emph{Polis}, pertença, dever cívico\\
\textbf{Conceito Grego:} \emph{Koinonia} (comunidade, partilha)

\part{Círculo III --- No Mundo (Cosmos)}

\chapter{Sobre o fracasso}\label{sobre-o-fracasso}

\emph{«Não desejes que as coisas aconteçam como queres. Deseja que
aconteçam como acontecem, e viverás bem.»}

--- Epicteto

\textbf{Âncora Estóica:} \emph{Amor Fati} (amor ao destino)\\
\textbf{Conceito Grego:} \emph{Telos} (propósito, fim)

\chapter{Vais morrer; vive agora}\label{vais-morrer-vive-agora}

\emph{«Podes sair da vida agora. Deixa que isso determine o que fazes,
dizes e pensas.»}

--- Marco Aurélio

\textbf{Âncora Estóica:} \emph{Memento Mori} (lembra-te que vais
morrer)\\
\textbf{Conceito Grego:} \emph{Kairos} (o momento certo, o agora)

\chapter{Sobre a paciência}\label{sobre-a-paciuxeancia}

\emph{«O tempo é um rio de acontecimentos que passam, e forte é a sua
corrente.»}

--- Marco Aurélio

\textbf{Âncora Estóica:} O tempo como mestre\\
\textbf{Conceito Grego:} \emph{Kairos} vs \emph{Chronos} (momento certo
vs tempo cronológico)

\chapter{Uma carta sem palavras}\label{uma-carta-sem-palavras}

\emph{«A posse que ninguém pode tirar ao homem é a paideia.»}

--- Menandro


\backmatter


\end{document}
